
\section{Additional Experiments}
\label{sec:additional-experiments}

We complement the main evaluation with experiments that broaden the
distribution of attack instances, control target difficulty across
architectures, and further examine robustness and computational cost.
Unless otherwise stated, we use ConvNetBN, the same poison budget $m$
and attack-specific optimization hyperparameters as in the main
experiments, and the default selector configuration
$\lambda=1$, $\alpha=2$, and $K=20$. The surrogate ensemble consists of
$20$ independently trained models instantiated from a fixed base seed,
which is also the ensemble reused by the main experiments.

\paragraph{Notation.} We write an attack instance as
$a\!\rightarrow\!b$, where $a$ is the ground-truth class of the target
test example and $b$ is the adversarial class from which the poison
bases are drawn. In this notation the two pairs of the main sweep,
labelled dog--bird and frog--airplane in
Table~\ref{tab:sweep_asr_clean_minus_cta}, are
bird$\rightarrow$dog and airplane$\rightarrow$frog; the labels there
list the adversarial class first. We use the arrow notation throughout
this appendix.

\paragraph{Protocol.} For every configuration, we select five target
test examples before running any attack and evaluate each target using
five independently initialized victims. Target examples are sampled
uniformly without replacement from the corresponding target class using
a fixed sampling seed. Following the main experiments, we exclude test
examples that the clean victim ensemble already assigns to the
adversarial class, since these are successes for any attack; no other
filtering is applied and no target is selected on the basis of an attack
outcome. Victims differ only in their initialization and batch ordering,
which are derived from the target index and a victim index; because
GPU kernels are not bitwise deterministic, matched seeds give matched
initializations rather than identical trained models. Random and DPP are
always evaluated using the same targets, victim indices, poison budget,
and downstream attack hyperparameters. The exact target indices are
reported in Appendix~\ref{app:additional-targets}.

These additional experiments use a reduced protocol relative to the main
sweep, which evaluates ten targets with six victims, in order to keep the
extended evaluation within our compute budget. Every comparison reported
in a given table uses the identical protocol, so the Random--DPP
contrast within a row is unaffected; absolute ASR values are not directly
comparable to the main tables.

\paragraph{Reporting.} We report ASR by first computing the success rate
of each target over its five victim runs and then averaging over the
five targets. With five targets and five victims, a single cell is
resolved to four percentage points and the sampling error over targets is
correspondingly large, so we treat individual cells as indicative and
base our conclusions on the direction and consistency of the
Random--DPP difference across configurations. Because Random and DPP are
evaluated on the same targets, we report the difference $\Delta$ as a
paired quantity and, where it is the quantity of interest, accompany it
with a $95\%$ bootstrap confidence interval obtained by resampling
targets.


\subsection{Generalization}
\label{sec:broader-generalization}

The main evaluation in Table~\ref{tab:sweep_asr_clean_minus_cta}
considers two attack instances, bird$\rightarrow$dog and
airplane$\rightarrow$frog. We conduct two additional evaluations to
determine whether the observed advantage of base selection extends
beyond them.

\paragraph{Additional CIFAR-10 class pairs.}
We extend the CIFAR-10 evaluation to six additional ordered class pairs:
$\mathrm{cat}\!\rightarrow\!\mathrm{dog}$,
$\mathrm{deer}\!\rightarrow\!\mathrm{horse}$,
$\mathrm{automobile}\!\rightarrow\!\mathrm{truck}$,
$\mathrm{bird}\!\rightarrow\!\mathrm{airplane}$,
$\mathrm{ship}\!\rightarrow\!\mathrm{frog}$, and
$\mathrm{truck}\!\rightarrow\!\mathrm{cat}$,
where the class on the left is the ground-truth class of the target and
the class on the right is the desired adversarial class. These pairs are
fixed before attack evaluation and cover both semantically similar and
dissimilar class combinations. For each pair, we sample five test
examples from the target class and keep these targets fixed across all
selection methods and poisoning objectives. Candidate bases are drawn
from the corresponding adversarial-class training examples, following
the threat model in Section~\ref{sec:preliminaries}.

We evaluate FC, GM, and SAPA using ConvNetBN at
$\epsilon=5\times10^{-3}$, with the same poison budget, attack
optimization settings, and victim-training procedure used in the main
experiments. ConvNetBN is used throughout this appendix so that the six
additional pairs, the cross-dataset evaluation, and the redundancy and
cost analyses share one architecture; the cross-architecture question is
addressed separately in Table~\ref{tab:matched-architecture}. We note
that ConvNetBN is also the architecture with the largest Random--DPP gap
in the main sweep, so the absolute margins reported here should not be
read as representative of all architectures. Each target is evaluated
over five independently trained victims, and Random and DPP use identical
targets and victim indices. For reference, the last row repeats the main
pair bird$\rightarrow$dog under this protocol, so that the six new pairs
can be compared against a familiar instance rather than against the
main sweep, which uses a different number of targets and victims.
Table~\ref{tab:broad-cifar} reports the resulting ASR for each ordered
class pair and the average across the six additional pairs.





\begin{table}[t]
\centering
\caption{
Generalization across six additional ordered CIFAR-10 class pairs.
The notation $a\!\rightarrow\!b$ denotes target class $a$ and
adversarial class $b$. Experiments use ConvNetBN,
$\epsilon=5\times10^{-3}$, five fixed targets per pair, and five victim
runs per target. Entries report mean ASR (\%) across targets.
}
\label{tab:broad-cifar}
\small
\setlength{\tabcolsep}{4.8pt}
\renewcommand{\arraystretch}{1.08}
\begin{tabular}{l|cc|cc|cc}
\toprule
&
\multicolumn{2}{c|}{\textbf{FC}} &
\multicolumn{2}{c|}{\textbf{GM}} &
\multicolumn{2}{c}{\textbf{SAPA}} \\
\textbf{Target $\rightarrow$ Adv.} &
\textbf{Random} & \textbf{DPP} &
\textbf{Random} & \textbf{DPP} &
\textbf{Random} & \textbf{DPP} \\
\midrule
cat $\rightarrow$ dog
    & 40.0 & 100.0 & 44.0 & 80.0 & 36.0 & 80.0 \\
deer $\rightarrow$ horse
    & 8.0 & 100.0 & 4.0 & 72.0 & 4.0 & 72.0 \\
automobile $\rightarrow$ truck
    & 4.0 & 68.0 & 12.0 & 48.0 & 12.0 & 48.0 \\
bird $\rightarrow$ airplane
    & 24.0 & 88.0 & 12.0 & 80.0 & 32.0 & 64.0 \\
ship $\rightarrow$ frog
    & 0.0 & 4.0 & 0.0 & 0.0 & 0.0 & 0.0 \\
truck $\rightarrow$ cat
    & 0.0 & 28.0 & 0.0 & 0.0 & 0.0 & 4.0 \\
\midrule
\textbf{Average}
    & 12.7 & 64.7 & 12.0 & 46.7 & 14.0 & 44.7 \\
\midrule
bird $\rightarrow$ dog \small{(main pair, this protocol)}
    & 24.0 & 80.0 & 12.0 & 64.0 & 24.0 & 68.0 \\
\bottomrule
\end{tabular}
\end{table}


\paragraph{Cross-dataset evaluation.}
We next evaluate whether the same selector remains effective outside
CIFAR-10, using CIFAR-100. Unlike the rest of this appendix we use
ResNet18 rather than ConvNetBN here: the $100$-way problem requires a
stronger backbone for the clean victim to be an informative base
classifier. ResNet18 is trained with the identical schedule and the
selector itself is unchanged. We sample five ordered
target--adversarial-class pairs uniformly without replacement from the
$100$ CIFAR-100 classes using a fixed seed, and then sample one target
test example for each pair. These five attack instances are fixed before
poison construction and shared between Random and DPP. Their class names
and target indices are listed in
Appendix~\ref{app:additional-targets}.

We additionally evaluate SVHN, which shares CIFAR-10's $32\times32$ inputs
and ten classes but consists of street-view house numbers rather than
natural objects. Because the input size and class count match CIFAR-10,
SVHN keeps ConvNetBN; the SVHN row therefore isolates a change of data
distribution, whereas the CIFAR-100 row changes the dataset and the
backbone together. SVHN's training split contains $73{,}257$ images, so
the same budget fraction buys $146$ poisons per target rather than
CIFAR-10's $100$ at $2\times10^{-3}$; the budget is a fraction of the
training set, and that is the quantity held fixed across rows. The five
digit pairs are sampled by the same seeded procedure, which for ten
classes consumes all of them exactly once.

We evaluate GM and SAPA at $\epsilon=5\times10^{-3}$ on CIFAR-10 and at
$\epsilon=2\times10^{-3}$ on CIFAR-100 and SVHN. The budgets differ because the
quantity that determines how much freedom the selector has is the ratio
between the poison budget and the size of the candidate pool, not the
budget alone. CIFAR-100 contains $500$ training images per class, one
tenth of the CIFAR-10 pool, so a budget of $1\times10^{-2}$ of the
$50{,}000$ training images would require $500$ poisons and consume the
entire adversarial class; every selection method would then return the
same set and the comparison would be vacuous by construction. At
$2\times10^{-3}$ the attacker selects $100$ of $500$ candidates. The
corresponding CIFAR-10 setting selects $250$ of $5{,}000$, so the
CIFAR-100 selector still operates on a substantially smaller pool and a
larger fraction of it; this makes the CIFAR-100 comparison the more
constrained of the two. Each attack instance is evaluated over the five
victim runs specified above. The comparison of interest is the
improvement from Random to DPP within each dataset rather than the
absolute ASR across datasets, which differ in number of classes and
candidate-pool size.

\begin{table}[t]
\centering
\caption{
Cross-dataset evaluation of base selection. The CIFAR-10 and SVHN rows
use ConvNetBN and the CIFAR-100 rows use ResNet18, in all cases with
five fixed attack instances and five victim runs per target.
Perturbation budgets are $\epsilon=5\times10^{-3}$ on CIFAR-10 and
$\epsilon=2\times10^{-3}$ on CIFAR-100 and SVHN, chosen so that the
poison budget does not exhaust the $500$-image candidate pool of a
CIFAR-100 class. Within each dataset, Random and DPP use identical
targets and attack settings. $\Delta$ denotes DPP minus Random.
}
\label{tab:cross-dataset}
\small
\setlength{\tabcolsep}{7pt}
\renewcommand{\arraystretch}{1.08}
\begin{tabular}{llccc}
\toprule
\textbf{Dataset} &
\textbf{Attack} &
\textbf{Random} &
\textbf{DPP} &
$\boldsymbol{\Delta}$ \\
\midrule
CIFAR-10 \small{(ConvNetBN)}     & GM   & 12.0 & 64.0 & +52.0 \\
CIFAR-10 \small{(ConvNetBN)}     & SAPA & 24.0 & 68.0 & +44.0 \\
\midrule
CIFAR-100 \small{(ResNet18)}      & GM   & 12.0 & 24.0 & +12.0 \\
CIFAR-100 \small{(ResNet18)}      & SAPA & 4.0 & 16.0 & +12.0 \\
\midrule
SVHN \small{(ConvNetBN)}          & GM   & 36.0 & 80.0 & +44.0 \\
SVHN \small{(ConvNetBN)}          & SAPA & 28.0 & 88.0 & +60.0 \\
\bottomrule
\end{tabular}
\end{table}


\paragraph{Matched-target architecture evaluation.}
The architecture results in the main sweep use independently sampled
target sets for different model--attack configurations. We therefore
perform an additional controlled comparison in which the attacked
examples are held fixed across architectures. We use the
bird$\rightarrow$dog instance of the main sweep, that is five bird-class
targets with dog as the adversarial class, and evaluate the exact same
targets with ConvNetBN, ResNet20BN, and VGG13. For every architecture,
the Random and DPP variants use the same five victim indices.

We evaluate GM and SAPA at
$\epsilon=5\times10^{-3}$ and retain the architecture-specific training
and attack hyperparameters used in the main experiments. Base selection
is performed using surrogates of the corresponding architecture, with
$K=20$, $\lambda=1$, and $\alpha=2$ in every case.
Table~\ref{tab:matched-architecture} therefore varies only the model
architecture while keeping the target set, adversarial class,
perturbation budget, attack objective, selector configuration, and
victim seeds matched.

\begin{table}[t]
\centering
\caption{
Matched-target evaluation across architectures. All configurations use
the same five bird$\rightarrow$dog targets,
$\epsilon=5\times10^{-3}$, and five victim runs per target.
Entries report mean ASR (\%) across targets.
}
\label{tab:matched-architecture}
\small
\setlength{\tabcolsep}{8pt}
\renewcommand{\arraystretch}{1.08}
\begin{tabular}{l|cc|cc}
\toprule
&
\multicolumn{2}{c|}{\textbf{GM}} &
\multicolumn{2}{c}{\textbf{SAPA}} \\
\textbf{Architecture} &
\textbf{Random} & \textbf{DPP} &
\textbf{Random} & \textbf{DPP} \\
\midrule
ConvNetBN  & 12.0 & 64.0 & 24.0 & 68.0 \\
ResNet20BN & 48.0 & 80.0 & 52.0 & 80.0 \\
VGG13      & 72.0 & 80.0 & 72.0 & 80.0 \\
\bottomrule
\end{tabular}
\end{table}


\subsection{Augmentation Robustness}
\label{sec:augmentation-aware}

Table~\ref{tab:augmentation-robustness} keeps the optimized poisons
fixed and changes only the victim-side augmentation. The poisons in that
table are constructed with the differentiable augmentation used by our
attacks throughout, which is not matched to the victim pipeline, so the
experiment measures robustness to an unanticipated change in victim
training. It does not distinguish sensitivity of the selected bases from
sensitivity of the downstream poison optimizer. We therefore additionally
evaluate an attacker that accounts for the victim transformation during
poison construction.

We use ConvNetBN on dog$\rightarrow$bird and the same five attack targets
throughout. We evaluate GM and SAPA at
$\epsilon=5\times10^{-3}$ under Crop+Flip and RandAugment. Crop+Flip
uses random $32\times32$ crops with four-pixel padding followed by a
horizontal flip with probability $0.5$. RandAugment uses the same
Crop+Flip pipeline followed by RandAugment with two operations and
magnitude $9$, matching the augmentation setting used in
Table~\ref{tab:augmentation-robustness}.

For each augmentation, we compare two poison-construction procedures.
In the \emph{unmatched} setting, the attack is optimized on the
unaugmented training examples and augmentation is introduced only when
training the victim. In the \emph{matched} setting, transformations from
the same augmentation distribution are sampled during poison
optimization and during victim training. All other attack settings,
including the poison budget, optimization iterations, perturbation
constraint, and victim-training schedule, are unchanged.

Matched optimization requires the augmentation to be differentiable, so
that transformed poisons can be backpropagated through during poison
construction. This holds for Crop+Flip, which we sample from the same
differentiable augmentation family used by the attack. RandAugment
includes discrete photometric operations that admit no such gradient, so
a matched attacker is not defined for it under our threat model; we
therefore report RandAugment in the unmatched setting only and mark the
matched entries as not applicable. Table~\ref{tab:augmentation-aware}
reports Random and DPP under both conditions.

\begin{table}[t]
\centering
\caption{
Augmentation-aware poison optimization on dog$\rightarrow$bird using
ConvNetBN and $\epsilon=5\times10^{-3}$. ``Unmatched'' applies
augmentation only during victim training; ``Matched'' additionally uses
the same augmentation distribution during poison optimization. Matched
optimization is defined only for differentiable augmentations, so it does
not apply to RandAugment. All settings use the same five targets and five
victim runs per target. Entries report mean ASR (\%).
}
\label{tab:augmentation-aware}
\small
\setlength{\tabcolsep}{5.5pt}
\renewcommand{\arraystretch}{1.08}
\begin{tabular}{ll|cc|cc}
\toprule
&
&
\multicolumn{2}{c|}{\textbf{Unmatched}} &
\multicolumn{2}{c}{\textbf{Matched}} \\
\textbf{Attack} &
\textbf{Augmentation} &
\textbf{Random} & \textbf{DPP} &
\textbf{Random} & \textbf{DPP} \\
\midrule
GM   & Crop+Flip   & 4.0 & 20.0 & 4.0 & 16.0 \\
GM   & RandAugment & 0.0 & 8.0 & \multicolumn{2}{c}{n/a} \\
\midrule
SAPA & Crop+Flip   & 0.0 & 16.0 & 0.0 & 16.0 \\
SAPA & RandAugment & 4.0 & 4.0 & \multicolumn{2}{c}{n/a} \\
\bottomrule
\end{tabular}
\end{table}


\subsection{Defense Evaluation}
\label{sec:utility-matched-defense}

The defense results in Table~\ref{tab:defense-robustness} use the
default configurations of EPIC and FRIENDS. Because these settings can
produce different reductions in clean accuracy, part of the observed
decrease in ASR may result from degradation of the classifier itself.
We therefore complement the default-setting evaluation with a
clean-utility-matched comparison.

We first train ConvNetBN on clean CIFAR-10 data while varying the
strength of each defense. Defense calibration uses no poisoned data and
no attack outcomes. For EPIC we vary the fraction of the training set
retained at each selection round, and for FRIENDS the magnitude of the
friendly noise; in both cases we order configurations from weakest to
strongest by that parameter and choose the strongest configuration whose
clean test accuracy remains within two percentage points of the
undefended ConvNetBN trained under the identical schedule. Once selected,
the defense hyperparameters are frozen and used identically for Random
and DPP.

Because the calibration reference is the undefended ConvNetBN under this
protocol, the comparison in Table~\ref{tab:utility-defense} is internal
to that table. Table~\ref{tab:defense-robustness} uses ResNet20BN with
the default defense configurations, so absolute ASR values should not be
compared across the two tables; the quantity of interest here is whether
the Random--DPP ordering survives once the defense is prevented from
recovering ASR by degrading the classifier.

We evaluate GM and SAPA on both bird$\rightarrow$dog and
airplane$\rightarrow$frog at $\epsilon=2\times10^{-2}$. For each
pair--attack combination, we use five fixed targets and the same five
victim seeds for Random and DPP. The poison budget, poison optimization,
and victim-training settings are otherwise identical to the
corresponding undefended experiments. Table~\ref{tab:utility-defense}
reports the clean accuracy obtained by each calibrated defense and the
resulting attack success.

\begin{table}[t]
\centering
\caption{
Utility-matched evaluation of EPIC and FRIENDS using ConvNetBN and
$\epsilon=2\times10^{-2}$. Defense strength is chosen using clean data
and subsequently fixed for all poisoned runs. Clean CTA is the accuracy
of the defended model trained without poisons; Random and DPP report
mean ASR (\%) across five targets and five victim runs per target.
}
\label{tab:utility-defense}
\small
\setlength{\tabcolsep}{5.5pt}
\renewcommand{\arraystretch}{1.08}
\begin{tabular}{lllccc}
\toprule
\textbf{Defense} &
\textbf{Target $\rightarrow$ Adv.} &
\textbf{Attack} &
\textbf{Clean CTA} &
\textbf{Random} &
\textbf{DPP} \\
\midrule
EPIC &
bird $\rightarrow$ dog &
GM &
80.33 & 36.0 & 52.0 \\

EPIC &
bird $\rightarrow$ dog &
SAPA &
80.33 & 36.0 & 64.0 \\

EPIC &
airplane $\rightarrow$ frog &
GM &
80.33 &  0.0 &  0.0 \\

EPIC &
airplane $\rightarrow$ frog &
SAPA &
80.33 &  0.0 &  0.0 \\

\midrule

FRIENDS &
bird $\rightarrow$ dog &
GM &
80.02 & 28.0 & 72.0 \\

FRIENDS &
bird $\rightarrow$ dog &
SAPA &
80.02 & 32.0 & 80.0 \\

FRIENDS &
airplane $\rightarrow$ frog &
GM &
80.02 &  0.0 &  8.0 \\

FRIENDS &
airplane $\rightarrow$ frog &
SAPA &
80.02 &  0.0 & 20.0 \\
\bottomrule
\end{tabular}
\end{table}


\subsection{Base Redundancy}
\label{sec:base-redundancy}

Greedy and DPP use the same pointwise score $s_i(x_t)$, but DPP
additionally discourages selecting candidates with similar
representations. We isolate this effect by controlling how many distinct
clean examples are used to initialize a poison set while holding the
total poison budget fixed.

This experiment uses ConvNetBN, SAPA, bird$\rightarrow$dog,
$\epsilon=5\times10^{-3}$, and the same five targets and victim indices
used above. Under our threat model a poison replaces a specific clean
training example, so the poison budget cannot be spent on repeated copies
of one image: $m$ copies of a single base occupy one training slot and
would reduce the effective budget from $m$ to one. We therefore vary
redundancy through the region a poison set occupies rather than through
literal duplication. Given the standard budget $m$, we take the $r$
highest-scoring candidates under $s_i(x_t)$ as seeds and fill the
remaining $m-r$ slots with the seeds' nearest neighbours in the surrogate
representation, cycling over the seeds so that the neighbourhoods grow at
the same rate. All $m$ selected bases are distinct and every one receives
its own independently optimized perturbation satisfying
$\|\delta_i\|_\infty\leq\epsilon$. The number of poisoned training
examples, the optimization objective, the number of optimization steps,
and the victim-training procedure are therefore identical across
configurations; only the number of distinct favourable regions the budget
is spread over varies, with $r\in\{1,2,5,10,m\}$.

Top-$1$ concentrates the entire budget in the single neighbourhood
around the highest-scoring base, while Top-$m$ places every poison at a
separately scored candidate and is equivalent to Greedy. Random and DPP
provide the two reference points. Comparing Top-$r$ against DPP separates
two effects that the pointwise score alone cannot: whether it is
sufficient to exploit the strongest region repeatedly, or whether
spreading the same budget over distinct regions is what matters.
Table~\ref{tab:base-redundancy} reports the resulting ASR.

\begin{table}[t]
\centering
\caption{
Effect of base redundancy at fixed poison budget $m$. Top-$r$ concentrates
the budget in the neighbourhoods of the $r$ highest-scoring clean
examples, using $m$ distinct and independently optimized poisons in every
configuration; Top-$m$ is equivalent to Greedy. Experiments use
ConvNetBN, SAPA, bird$\rightarrow$dog, and
$\epsilon=5\times10^{-3}$. Entries report mean ASR (\%) across five
targets and five victim runs per target.
}
\label{tab:base-redundancy}
\small
\setlength{\tabcolsep}{6pt}
\renewcommand{\arraystretch}{1.08}
\begin{tabular}{c|ccccccc}
\toprule
&
\textbf{Random} &
\textbf{Top-1} &
\textbf{Top-2} &
\textbf{Top-5} &
\textbf{Top-10} &
\textbf{Greedy} &
\textbf{DPP} \\
\midrule
\textbf{ASR (\%)} &
 24.0 & 52.0 & 52.0 & 32.0 & 48.0 & 60.0 & 68.0 \\
\bottomrule
\end{tabular}
\end{table}


\subsection{Computational Cost}
\label{sec:computational-cost}

We finally quantify the computation introduced by base selection.
Measurements use ConvNetBN, the default $K=20$ surrogate ensemble, and
the full adversarial-class candidate pool used in the CIFAR-10
experiments. All timings are collected on the same hardware and averaged
over the five bird$\rightarrow$dog targets.

The selector has two types of cost. Training the $20$ surrogate models
is performed once on the clean training set and the resulting models can
be reused across targets and poisoning objectives. For each new target,
we compute $B_i$ and $R_i(x_t)$ for every adversarial-class candidate,
form the corresponding candidate scores, and then perform either Greedy
or DPP subset selection. Greedy and DPP therefore share the surrogate
training and candidate-scoring cost, while only DPP incurs the
diversity-aware subset-selection step.

The upper part of Table~\ref{tab:computational-cost} reports wall-clock
time and peak GPU memory for these components. The lower part reports
the complete attack-construction time for FC, GM, and SAPA. For each
attack, Random and DPP use the same target, poison budget, optimization
iterations, and attack hyperparameters. DPP runtime includes all
target-dependent selection computation but excludes the one-time
surrogate-training cost, which is reported separately above. The
overhead column reports the relative increase in total per-target
attack-construction time from replacing Random with DPP selection. This
ratio is attack-dependent by construction: the selection cost is the same
for all three attacks, so it is largest relative to FC, whose poison
optimization is by far the cheapest, and smallest relative to GM and
SAPA. The absolute selection time is the quantity to compare across
attacks.

\begin{table}[t]
\centering
\caption{
Computational cost of base selection and poison construction using
ConvNetBN. Surrogate training is performed once and reused across
targets and poisoning objectives; all remaining timings are per target.
Measurements use the same hardware and are averaged over five targets.
}
\label{tab:computational-cost}
\small
\setlength{\tabcolsep}{6pt}
\renewcommand{\arraystretch}{1.08}
\begin{tabular}{lccc}
\toprule
\multicolumn{4}{c}{\textbf{Selection cost}} \\
\midrule
\textbf{Stage} &
\textbf{Time} &
\textbf{Peak Memory} &
\textbf{Frequency} \\
\midrule
Train $K=20$ surrogates & 1220 s & 1748.2 MB & Once \\
Candidate scoring       & 0.68 s & 1336.7 MB & Per target \\
Greedy subset selection & $<0.01$ s & 1336.7 MB & Per target \\
DPP subset selection    & 0.19 s & 3208.6 MB & Per target \\
\midrule
\multicolumn{4}{c}{\textbf{End-to-end attack cost}} \\
\midrule
\textbf{Attack} &
\textbf{Random} &
\textbf{DPP} &
\textbf{Overhead (\%)} \\
\midrule
FC   & 5.0 s & 5.9 s & 17.3 \\
GM   & 241.0 s & 241.9 s & 0.4 \\
SAPA & 242.0 s & 242.9 s & 0.4 \\
\bottomrule
\end{tabular}
\end{table}


\subsection{Additional Target Sets}
\label{app:additional-targets}

For reproducibility, we report the exact target examples used in the
additional experiments after sampling them according to the protocol
above. CIFAR-10 indices refer to the standard test split. CIFAR-100
entries report the sampled class pair and the index of the target within
the standard CIFAR-100 test split.

\begin{table}[t]
\centering
\caption{
Target sets used in the additional experiments. All targets are fixed
before attack construction and shared between Random and DPP.
}
\label{tab:additional-target-indices}
\small
\setlength{\tabcolsep}{5pt}
\renewcommand{\arraystretch}{1.08}
\begin{tabular}{lll}
\toprule
\textbf{Dataset} &
\textbf{Target $\rightarrow$ Adv.} &
\textbf{Target indices} \\
\midrule
CIFAR-10 & cat $\rightarrow$ dog
    & 2034, 607, 4375, 5493, 5530 \\
CIFAR-10 & deer $\rightarrow$ horse
    & 2180, 1341, 5296, 2353, 5414 \\
CIFAR-10 & automobile $\rightarrow$ truck
    & 2477, 8215, 726, 261, 1396 \\
CIFAR-10 & bird $\rightarrow$ airplane
    & 1083, 160, 6277, 9915, 8956 \\
CIFAR-10 & ship $\rightarrow$ frog
    & 1347, 7485, 3695, 132, 73 \\
CIFAR-10 & truck $\rightarrow$ cat
    & 9378, 6579, 4761, 2174, 9743 \\
CIFAR-10 & bird $\rightarrow$ dog
    & 3741, 6857, 8252, 1126, 3522 \\
\midrule
CIFAR-100 & willow\_tree $\rightarrow$ sea
    & 3818 \\
CIFAR-100 & bicycle $\rightarrow$ plain
    & 4646 \\
CIFAR-100 & lawn\_mower $\rightarrow$ wardrobe
    & 3576 \\
CIFAR-100 & road $\rightarrow$ bottle
    & 4167 \\
CIFAR-100 & cattle $\rightarrow$ sunflower
    & 4276 \\
\midrule
SVHN & 2 $\rightarrow$ 9 & 14586 \\
SVHN & 1 $\rightarrow$ 6 & 23802 \\
SVHN & 3 $\rightarrow$ 8 & 11465 \\
SVHN & 5 $\rightarrow$ 7 & 59 \\
SVHN & 4 $\rightarrow$ 0 & 7975 \\
\bottomrule
\end{tabular}
\end{table}